% answers_short.tex —— 2011 年·全国新课标卷·理科 · 简短答案速查
% 由 tools/make_short_answers.py 自 answers.tex 的 \gksolhead/\gkanswer 提取，
% 并与 2011/2011_数学真题_全国新课标理_答案要点.md 逐题比对后生成，勿手改（改源 MD 后重跑）。

\begin{gknotice}
\textbf{使用说明}：本节只给最终答案便于快速核对，完整推导见解析版。答案为编写者独立推导，\textbf{非官方标准答案}。
\end{gknotice}

\gkpart{选择题}
\begin{center}\begin{tabular}{ccccccccccccc}
\toprule
题号 & 1 & 2 & 3 & 4 & 5 & 6 & 7 & 8 & 9 & 10 & 11 & 12 \\
\midrule
答案 & C & B & A & A & B & D & B & D & C & A & A & D \\
\bottomrule
\end{tabular}\end{center}

\gkpart{填空题}
\begin{itemize}[leftmargin=2.6em]
  \item[\textbf{13.}] $-6$
  \item[\textbf{14.}] $\dfrac{x^2}{16}+\dfrac{y^2}{8}=1$
  \item[\textbf{15.}] $8\sqrt3$
  \item[\textbf{16.}] $2\sqrt7$
\end{itemize}

\gkpart{解答题}
\begin{itemize}[leftmargin=2.6em]
  \item[\textbf{17.}] (1) $a_n=3^{-n}$；(2) $b_n=-\dfrac{n(n+1)}2$，$\displaystyle\sum_{k=1}^n\dfrac1{b_k}=-\dfrac{2n}{n+1}$
  \item[\textbf{18.}] (1) 见解析版；(2) 二面角 $A-PB-C$ 的余弦值为 $-\dfrac{2\sqrt7}{7}$
  \item[\textbf{19.}] (1) A 配方优质品率 $30\%$，B 配方优质品率 $42\%$；(2) $P(X=-2)=0.04$，$P(X=2)=0.54$，$P(X=4)=0.42$，$E(X)=2.68=\dfrac{67}{25}$
  \item[\textbf{20.}] (1) $C:y=\dfrac{x^2}{4}-2$；(2) 原点到切线的距离最小值为 $2$
  \item[\textbf{21.}] (1) $a=b=1$；(2) $k\in(-\infty,0]$
  \item[\textbf{22.}] (1) 见解析版；(2) 所在圆的半径为 $5\sqrt2$
  \item[\textbf{23.}] (1) $C_2:x^2+(y-4)^2=16$；(2) $|AB|=2\sqrt3$
  \item[\textbf{24.}] (1) $(-\infty,-1]\cup[3,+\infty)$；(2) $a=2$
\end{itemize}
